% !TEX TS-program = xelatex
% !TEX encoding = UTF-8 Unicode
% !Mode:: "TeX:UTF-8"

\documentclass{resume}
\usepackage{zh_CN-Adobefonts_external} % Simplified Chinese Support using external fonts (./fonts/zh_CN-Adobe/)
% \usepackage{NotoSansSC_external}
% \usepackage{NotoSerifCJKsc_external}
% \usepackage{zh_CN-Adobefonts_internal} % Simplified Chinese Support using system fonts
\usepackage{linespacing_fix} % disable extra space before next section
\usepackage{cite}

\begin{document}
\pagenumbering{gobble} % suppress displaying page number

\name{刘冠云}

\basicInfo{
\email{g.liu1@ufl.edu} \textperiodcentered\ 
\phone{(+1) 352-278-5052} \textperiodcentered\ 
  % \linkedin[billryan8]{https://www.linkedin.com/in/billryan8}
  \home{3611 SW 34th ST, Apt. 240, Gainesville, FL-32608}
}
 
\section{\faGraduationCap\  教育背景}
\datedsubsection{\textbf{佛罗里达大学}，佛罗里达，美国}{2023 -- 至今}
\textit{在读博士研究生}\ 机械与航空航天工程，预计 2027 年 8 月毕业
\datedsubsection{\textbf{约翰霍普金斯大学}，马里兰，美国}{2021 -- 2023}
\textit{理学硕士}\ 机械工程 
\datedsubsection{\textbf{帝国理工学院}，伦敦，英国}{2017 -- 2018}
\textit{理学硕士$\cdot$卓越学位(Distinction)}\ 高级航空工程 
\datedsubsection{\textbf{布里斯托大学}，布里斯托，英国}{2014 -- 2017}
\textit{工学学士$\cdot$一等学位(First Class Honours)}\ 航空航天工程 

\section{\faFile\ 学术论文}
\subsection{期刊论文：}
\begin{itemize}
  \item \textbf{Liu, G.}, \& Menezes, A. A. (2024). Control of a Noncooperative Positive Nonlinear System by Augmented Positive Linear System Regulation. IEEE Control Systems Letters.
  \item Schwertheim, A., Rosati Azevedo, E., \textbf{Liu, G.}, Bosch Borràs, E., Bianchi, L., \& Knoll, A. (2021). Interlaboratory validation of a hanging pendulum thrust balance for electric propulsion testing. Review of Scientific Instruments, 92(3).
\end{itemize}
\subsection{会议论文：}
\begin{itemize}
  \item Wang, Y., Al-Zogbi, L., \textbf{Liu, G.}, Liu, J., Tokuda, J., Krieger, A., \& Iordachita, I. (2024, May). Bevel-tip needle deflection modeling, simulation, and validation in multi-layer tissues. In 2024 IEEE International Conference on Robotics and Automation (ICRA) (pp. 11598-11604). IEEE.
  \item \textbf{Liu, G.}, Wang, Y., Li, G., Cleary, K., \& Iordachita, I. (2022, October). Evaluation of Needle Driver Designs for Robot-Assisted Needle Insertions Under MRI Guidance. In ASME International Mechanical Engineering Congress and Exposition (Vol. 86663, p. V004T06A020). American Society of Mechanical Engineers.
  \item Wang, Y., \textbf{Liu, G.}, Li, G., Cleary, K., \& Iordachita, I. (2022, July). An MR-Conditional Needle Driver for Robot-Assisted Spinal Injections: Design Modifications and Evaluations. In 2022 44th Annual International Conference of the IEEE Engineering in Medicine \& Biology Society (EMBC) (pp. 3307-3312). IEEE.
\end{itemize}

\section{\faList\ 科研项目}
\datedsubsection{\textbf{院前战场救治中个性化自动止血控制的研究}}{2024年4月 -- 至今}
\role{数据分析，数学建模}{DARPA 项目}
\vspace{-0.7em}
\begin{onehalfspacing}
\begin{itemize}
  \item 分析18种凝血因子对凝血酶生成和血凝块强度的控制作用
  \item 研究凝血因子浓度随时间的动态变化规律
  \item 开发时变模型，为优化凝血反应的控制系统提供支持
\end{itemize}
\end{onehalfspacing}

\datedsubsection{\textbf{凝血酶生成控制理论开发以治疗凝血功能障碍}}{2023年8月 -- 至今}
\role{理论研究，算法开发}{美国国家科学基金会项目}
\vspace{-0.7em}
\begin{onehalfspacing}
  \begin{itemize}
    \item 根据CAT和TEG实验数据创建用于预测CAT和TEG诊断指标的传递函数数学模型
    \item 基于CAT和TEG的传递函数模型并结合生物信息学的相关生化反应，构建相应的正向非线性状态空间模型
    \item 基于传统的非线性系统的控制器设计理论，开发可用于正向非线性系统的控制器设计理论
    \item 使用MATLAB，SimuLink和Python验证所开发的控制器
  \end{itemize}
\end{onehalfspacing}

\datedsubsection{\textbf{创伤护理中氨甲环酸个性化给药的定量模型研究}}{2024年1月 -- 2024年9月}
\role{数据分析，数学建模}{美国国家科学基金会项目}
\vspace{-0.7em}
\begin{onehalfspacing}
  \begin{itemize}
    \item 设计并更新已有模型，根据凝血因子浓度和预期的TEG曲线预测创伤患者的个性化氨甲环酸（TXA）剂量
    \item 量化TXA对血液凝固的影响，实验证明TXA可抵消tPA的影响，延迟血凝块溶解，并稳定血凝块强度
    \item 将抗纤溶参数整合到模型中，建立TXA浓度与模型参数之间的双射映射，实现动态和精确的剂量推荐
    \item 开创性地分析合成患者血样，人工创建创伤患者样本，模拟凝血动力学并验证模型预测
  \end{itemize}
\end{onehalfspacing}

\datedsubsection{\textbf{非协同(Noncooperative)正向非线性系统的控制理论研究}}{2023年8月 -- 2024年9月}
\role{理论研究}{美国国家科学基金会项目}
\vspace{-0.7em}
\begin{onehalfspacing}
  \begin{itemize}
    \item 设计前馈和反馈控制器，解决生物学背景（如血液凝固）中非协同系统的挑战
    \item 引入附加控制器，恢复线性化系统的正向性，并确保其与非线性动态一致
    \item 采用混合建模方法，结合机理和现象学模型，简化控制器设计，同时保持生物学意义
    \item 提供正向性丧失的见解，引入鲁棒控制律，并提出生物医学及其他非线性系统的未来应用方向
  \end{itemize}
\end{onehalfspacing}

\datedsubsection{\textbf{使用SOFA进行医用针模型构建}}{2022 年5月 -- 2023年7月}
\role{C++编程与仿真}{个人项目}
\vspace{-0.7em}
\begin{onehalfspacing}
  \begin{itemize}
    \item 设计并制造用于实验的仪器
    \item 使用SOFA模拟医用针插入凝胶中的形变并收集模拟形变数据
    \item 通过CT扫描将针插入凝胶或肉中进行实验，收集针偏转数据
    \item 使用有限元理论并通过对比收集的模拟数据和实验数据构建在凝胶中医用针的形变模型
  \end{itemize}
\end{onehalfspacing}

\datedsubsection{\textbf{稳手眼机器人（SHER）的优化控制与RCM控制}}{2022 年8月 -- 2022年12月}
\role{最优化控制}{个人项目}
\vspace{-0.7em}
\begin{onehalfspacing}
  \begin{itemize}
    \item 推导了RCM约束的惩罚项表达式并将其融入成本函数
    \item 使用DDP算法找到将针尖移动到目标位置的最优轨迹
    \item 调整成本函数中的权重，使得针尖与目标位置的最大误差小于0.33毫米
  \end{itemize}
\end{onehalfspacing}

\datedsubsection{\textbf{核磁共振引导下的腰椎注射机器人}}{2021 年8月 -- 2022年5月}
\role{机械设计}{个人项目}
\vspace{-0.7em}
\begin{onehalfspacing}
  \begin{itemize}
    \item 设计了一个环锁机构，使临床医生能够快速安装针驱动装置
    \item 与之前的设计相比，在安装针驱动装置时，减少了高达67\%的力和73\%的扭矩
  \end{itemize}
\end{onehalfspacing}

\datedsubsection{\textbf{针对毫牛级等离子体推力器的推力天平开发与测试}}{2017 年8月 -- 2018年9月}
\role{机械设计}{硕士毕业设计项目}
\vspace{-0.7em}
\begin{onehalfspacing}
  \begin{itemize}
    \item 设计并制造了一种能够测量毫牛级推力的设备，用于等离子体推力器性能表征。
    \item 设计主动热控系统以防止传感器过热
    \item 将校准因子的误差降低约7\%
  \end{itemize}
\end{onehalfspacing}


% Reference Test
%\datedsubsection{\textbf{Paper Title\cite{zaharia2012resilient}}}{May. 2015}
%An xxx optimized for xxx\cite{verma2015large}
%\begin{itemize}
%  \item main contribution
%\end{itemize}

\section{\faCogs\ 专业技能}
% increase linespacing [parsep=0.5ex]
\begin{itemize}[parsep=0.5ex]
  \item 编程语言: Python，C++，MATLAB，\LaTeX
  \item 平台: Arduino，Linux，树莓派
  \item 软件: ROS (Robot Operating System)，SimuLink，SolidWorks，Autodesk Inventor，Mathematica，Vim，Emacs，Git，Illuminator，PhotoShop，Inkscape
\end{itemize}

% \section{\faUsers\ 课外经历}
% \datedline{\textit{第一名}, xxx 比赛}{2013 年6 月}
% \datedline{其他奖项}{2015}

% \section{\faInfo\ 其他}
% % increase linespacing [parsep=0.5ex]
% \begin{itemize}[parsep=0.5ex]
%   \item 技术博客: http://blog.yours.me
%   \item GitHub: https://github.com/username
%   \item 语言: 英语 - 熟练(TOEFL xxx)
% \end{itemize}

%% Reference
%\newpage
%\bibliographystyle{IEEETran}
%\bibliography{mycite}
\end{document}
